\section{Literature Review}
\label{sec:literature}

\subsection{Existing Research}

Other researchers have already studied parts of this problem, and their work gives this thesis a strong starting point.

First, research has shown that AI models can be taught to use tools on their own, without being told exactly when to use them. Schick et al.\ \cite{schick2023toolformer} demonstrated this with a system called Toolformer, which learned to call external tools by generating the call inside its own output. This established that tool use is possible, but the study only tested small numbers of tools called one at a time. It did not ask what happens when many tools are available at once.

Second, Patil et al.\ \cite{patil2023gorilla} tested exactly that question. They gave AI models access to over 1,600 different tools at once and measured how often the model picked the right one. They found that accuracy dropped significantly as the number of tools grew. Their benchmark, called Gorilla, is one of the first papers to directly show that giving an AI too many tools at once causes problems. This is the core problem that progressive disclosure is designed to solve.

Third, Qin et al.\ \cite{qin2023toolllm} extended this work to over 16,000 real-world tools (called ToolBench) and found that even the best models fail at multi-step tasks. This tells us that the problem is not just about picking the right tool --- it is also about following through on a series of steps correctly. This directly motivates the need for structured skill files with step-by-step execution policies.

Fourth, Yao et al.\ \cite{yao2022react} introduced a framework called ReAct, which makes AI agents write out their reasoning before taking each action. This produces a full record of every step the agent took, which makes it possible to see exactly where the agent went right or wrong. This study uses ReAct-style traces as the basis for measuring whether agents follow skill instructions correctly.

Fifth, Liu et al.\ \cite{liu2023agentbench} built a benchmark specifically for testing AI agents, called AgentBench. Their key finding was that final output is not enough to judge how well an agent works. An agent might get the right answer by accident, through a very messy process. They showed that the path the agent takes --- called the \textit{trajectory} --- must also be evaluated. This is why the current study measures both the final task result and the trajectory quality separately.

Sixth, Zhou et al.\ \cite{zhou2023webarena} tested frontier AI models on realistic, long multi-step tasks using a web browser. They found that even the best models succeed on fewer than 15 out of every 100 tasks when the full task path is evaluated. This confirms that measuring only the final output gives a misleadingly positive picture of agent performance.

Seventh, Shi et al.\ \cite{shi2024context} showed that when extra, irrelevant information is added to an AI model's input, its reasoning gets worse --- even if the model is a state-of-the-art system. They tested this by injecting distractors into math and reading tasks. The result directly explains why a long list of tool definitions should hurt agent performance: all the irrelevant tool descriptions act as distractors inside the model's context.

Finally, Liu et al.\ \cite{liu2024lost} discovered that AI models are bad at paying attention to information in the middle of a long context. Information near the beginning or end of the input is used reliably, but information in the middle is often ignored. This means that even if the right tool definition is present in a long flat list, the model may fail to use it correctly simply because it appeared in the middle of the list.

\subsection{Summary of Key Related Work (2022--2025)}

Table~\ref{tab:related-work} gives a quick overview of the most directly related studies, showing what each one found and what question it left unanswered.

\begin{table}[h!]
\centering
\small
\caption{Summary of Related Work (2022--2025)}
\label{tab:related-work}
\begin{tabular}{|p{0.65cm}|p{2.3cm}|p{3.5cm}|p{3.0cm}|p{3.9cm}|}
\hline
\textbf{Year} & \textbf{Study} & \textbf{Key Finding} & \textbf{Methodology} & \textbf{Gap Left} \\
\hline
2022 &
Yao et al.\ \cite{yao2022react} (ReAct) &
Writing out reasoning before each action improves task completion and produces readable traces &
Chain-of-thought prompting with tool call execution &
Does not compare flat tool lists against progressive disclosure \\
\hline
2023 &
Schick et al.\ \cite{schick2023toolformer} (Toolformer) &
AI models can learn to call tools by themselves without being told exactly when &
Fine-tuning on automatically generated API call examples &
Only tests single tool calls; no structured skill policies; no progressive loading \\
\hline
2023 &
Patil et al.\ \cite{patil2023gorilla} &
Tool selection accuracy falls as the number of available tools increases &
Benchmark with 1,600+ APIs comparing different models &
No procedural constraints; no trajectory evaluation; no skill hierarchy \\
\hline
2023 &
Qin et al.\ \cite{qin2023toolllm} (ToolBench) &
Even top models fail on multi-step tasks with large tool sets &
Benchmark with 16,000+ real-world tools across multiple steps &
No comparison of progressive loading vs.\ flat lists \\
\hline
2023 &
Liu et al.\ \cite{liu2023agentbench} (AgentBench) &
Final output alone is not enough to evaluate agents; the execution path matters too &
Eight interactive test environments with dual outcome and trajectory scoring &
Does not compare flat tool use against progressive disclosure \\
\hline
2023 &
Zhou et al.\ \cite{zhou2023webarena} (WebArena) &
Top AI models succeed on fewer than 15\% of long multi-step tasks when the full path is scored &
Realistic browser environment with full task evaluation &
Only tests a single tool-loading pattern \\
\hline
2024 &
Shi et al.\ \cite{shi2024context} &
Extra irrelevant information in an AI's input makes its reasoning worse &
Controlled tests adding distractors to math, reading, and coding tasks &
Not an agent study; does not test the effect of large tool lists specifically \\
\hline
2024 &
Liu et al.\ \cite{liu2024lost} &
AI models ignore information in the middle of long inputs &
Controlled experiments varying where key information appears in the context &
Not an agent study; no progressive loading comparison \\
\hline
2025 &
OpenClaw \cite{openclaw_2025} &
Three-tier skill loading claims to improve efficiency over flat tool lists &
Open-source autonomous agent framework with \texttt{SKILL.md} architecture &
No independent test; no controlled comparison against flat tool use \\
\hline
\end{tabular}
\end{table}

\subsection{Research Gap}

Even though all of these areas have been studied separately, nobody has put them together in one place. No study has directly compared an agent using a flat tool list against an agent using progressive skill disclosure on the same tasks, with the same models, measuring both the final result and the quality of the execution path. The OpenClaw framework claims that progressive disclosure works better, but that claim has never been independently tested. This thesis fills that gap.

\subsection{Research Questions}

Based on the work above, this thesis will answer three specific questions:

\begin{itemize}
    \item \textbf{Question 1:} Does an agent using progressive skill disclosure make significantly fewer hallucinated tool calls and follow task instructions more correctly than an agent using a flat tool list?
    \item \textbf{Question 2:} How many fewer tokens does progressive disclosure use per task compared to flat tool use, and does this difference grow as more tools are added to the library?
    \item \textbf{Question 3:} Does progressive disclosure help small local models more than it helps large cloud models, given that small models have shorter context windows and are more sensitive to long inputs?
\end{itemize}
